% ============================================================
\chapter{Produits D\'eriv\'es Exotiques}
\label{ch:exotiques}
% ============================================================

\section{Introduction}

Les options vanilles --- calls et puts europ\'eens --- sont les briques \'el\'ementaires de la finance des d\'eriv\'es. Mais les besoins des entreprises et des investisseurs sont infiniment plus vari\'es~: une compagnie a\'erienne veut se couvrir contre la \emph{moyenne} du prix du k\'eros\`ene sur un trimestre (option asiatique)~; un exportateur veut une protection qui s'active uniquement si le taux de change franchit un seuil (option barri\`ere)~; un investisseur veut le meilleur rendement parmi plusieurs actifs (option sur panier). Ces \emph{options exotiques}, n\'ees dans les salles de march\'e des ann\'ees 1980--1990, posent des d\'efis math\'ematiques consid\'erables~: la plupart n'admettent pas de formule ferm\'ee \`a la Black--Scholes et n\'ecessitent des m\'ethodes num\'eriques (Monte Carlo, diff\'erences finies, arbres).

\begin{intuition}
Les options vanilles sont des briques \'el\'ementaires. Les options exotiques
combinent ces briques ou ajoutent des conditions (barri\`eres, moyennes,
extr\'ema) pour cr\'eer des profils de gain sur mesure. Leur valorisation
requiert souvent des m\'ethodes num\'eriques avanc\'ees car les formules
ferm\'ees sont rares.
\end{intuition}

\section{Options \`a barri\`ere}

\begin{definition}[Option \`a barri\`ere]
\index{option \`a barri\`ere}
Une \textbf{option \`a barri\`ere} est une option dont l'existence ou
l'activation d\'epend du franchissement d'un niveau $B$ par le sous-jacent
$S_t$ durant la vie de l'option. On distingue :
\begin{itemize}
    \item \textbf{Knock-out} (d\'esactivante) : l'option dispara\^it si $S_t$
    atteint $B$.
    \item \textbf{Knock-in} (activante) : l'option n'existe que si $S_t$
    atteint $B$.
\end{itemize}
Combin\'e avec la direction (up/down), cela donne 8 types fondamentaux.
\end{definition}

\begin{theoreme}[Relation in-out]
\index{relation in-out}
Pour toute barri\`ere $B$, le prix d'un call (ou put) vanille v\'erifie :
\[
    V_{\mathrm{vanille}} = V_{\mathrm{knock\text{-}in}} + V_{\mathrm{knock\text{-}out}}.
\]
\end{theoreme}

\begin{theoreme}[Formule de Merton pour un down-and-out call]
\index{formule de Merton}
Dans le mod\`ele de Black-Scholes, pour un call down-and-out avec barri\`ere
$B < S_0$ et $B < K$ :
\[
    C_{\mathrm{do}} = C_{\mathrm{BS}} - \left(\frac{B}{S_0}\right)^{2\lambda}
    C_{\mathrm{BS}}\!\left(\frac{B^2}{S_0}, K, r, \sigma, T\right),
\]
o\`u $\lambda = \frac{r - \sigma^2/2}{\sigma^2} + \frac{1}{2}$ et
$C_{\mathrm{BS}}$ est le prix de Black-Scholes.
\end{theoreme}

\begin{attention}[title=\textbf{Risque de couverture des barri\`eres}]
Le delta d'une option barri\`ere diverge lorsque le sous-jacent approche de la
barri\`ere, rendant la couverture dynamique extr\^emement co\^uteuse. En
pratique, on utilise des barri\`eres discr\`etes (observ\'ees quotidiennement)
avec une correction de Broadie, Glasserman et Kou.
\end{attention}

\section{Options asiatiques}

\begin{definition}[Option asiatique]
\index{option asiatique}
Une \textbf{option asiatique} a un payoff d\'ependant de la moyenne du
sous-jacent :
\begin{itemize}
    \item \textbf{Average price} : payoff $= \left(\bar{S} - K\right)^+$
    (call) o\`u $\bar{S} = \frac{1}{T}\int_0^T S_t\,\mathrm{d}t$ (moyenne
    continue) ou $\bar{S} = \frac{1}{n}\sum_{i=1}^n S_{t_i}$ (moyenne
    discr\`ete).
    \item \textbf{Average strike} : payoff $= \left(S_T - \bar{S}\right)^+$.
\end{itemize}
\end{definition}

\begin{proposition}[Propri\'et\'es de l'asiatique]
\begin{enumerate}
    \item La moyenne r\'eduit la volatilit\'e effective :
    $\Var(\bar{S}) < \Var(S_T)$, donc les options asiatiques sont moins
    ch\`eres que les vanilles correspondantes.
    \item La distribution de $\bar{S}$ n'est pas lognormale, m\^eme si $S_t$
    suit un mouvement brownien g\'eom\'etrique.
    \item Il n'existe pas de formule ferm\'ee exacte dans Black-Scholes pour
    la moyenne arithm\'etique (contrairement \`a la moyenne g\'eom\'etrique).
\end{enumerate}
\end{proposition}

\begin{theoreme}[Prix de l'asiatique g\'eom\'etrique]
Pour la moyenne g\'eom\'etrique continue
$\bar{S}_G = \exp\!\left(\frac{1}{T}\int_0^T \ln S_t\,\mathrm{d}t\right)$,
le prix du call est donn\'e par la formule de Black-Scholes modifi\'ee avec
\[
    \hat{\sigma} = \frac{\sigma}{\sqrt{3}}, \qquad
    \hat{r} = \frac{1}{2}\left(r - \frac{\sigma^2}{6}\right).
\]
\end{theoreme}

\section{Options lookback}

\begin{definition}[Option lookback]
\index{option lookback}
Une \textbf{option lookback} d\'epend de l'extr\'emum du sous-jacent :
\begin{itemize}
    \item \textbf{Floating strike lookback call} :
    payoff $= S_T - \min_{0 \leq t \leq T} S_t$.
    \item \textbf{Fixed strike lookback call} :
    payoff $= \left(\max_{0 \leq t \leq T} S_t - K\right)^+$.
\end{itemize}
\end{definition}

\begin{theoreme}[Goldman, Sosin, Gatto (1979)]
\index{formule GSG}
Dans Black-Scholes, le prix du floating strike lookback call est
\[
    V = S_0\,\Phi(d_1) - S_0\,e^{-rT}\frac{\sigma^2}{2r}\,\Phi(-d_1)
    - S_{\min}\,e^{-rT}\left[\Phi(d_2) - \frac{\sigma^2}{2r}\left(\frac{S_0}{S_{\min}}\right)^{-2r/\sigma^2}\Phi(d_3)\right],
\]
o\`u $d_1, d_2, d_3$ sont des fonctions explicites de $S_0, S_{\min}, r, \sigma, T$.
\end{theoreme}

\section{Options binaires (digitales)}

\begin{definition}[Option binaire]
\index{option binaire}
Une \textbf{option binaire} (ou digitale) a un payoff discontinu :
\begin{itemize}
    \item \textbf{Cash-or-nothing call} : payoff $= Q \cdot \ind_{S_T > K}$
    (paye $Q$ si $S_T > K$).
    \item \textbf{Asset-or-nothing call} : payoff $= S_T \cdot \ind_{S_T > K}$.
\end{itemize}
\end{definition}

\begin{proposition}[Prix dans Black-Scholes]
\begin{align*}
    V_{\mathrm{cash}} &= Q\,e^{-rT}\,\Phi(d_2), \\
    V_{\mathrm{asset}} &= S_0\,\Phi(d_1),
\end{align*}
o\`u $d_1, d_2$ sont les quantit\'es standard de Black-Scholes.
Un call vanille se d\'ecompose : $C = V_{\mathrm{asset}} - K\,e^{-rT}\Phi(d_2)$.
\end{proposition}

\begin{remarque}
Le delta d'une option binaire diverge pr\`es du strike \`a l'\'ech\'eance,
ce qui rend la couverture tr\`es difficile. En pratique, on les r\'eplique par
un \textit{call spread} serr\'e.
\end{remarque}

\section{Valorisation par Monte Carlo}

\begin{definition}[Estimateur Monte Carlo]
\index{Monte Carlo!valorisation}
Le prix d'une option exotique de payoff $h(S_{t_1}, \ldots, S_{t_n})$ est
\[
    V_0 = e^{-rT}\,\E^{\Q}[h(S_{t_1}, \ldots, S_{t_n})]
    \approx \frac{e^{-rT}}{M}\sum_{j=1}^{M} h(S_{t_1}^{(j)}, \ldots, S_{t_n}^{(j)}),
\]
o\`u les trajectoires $(S^{(j)})$ sont simul\'ees sous la mesure risque-neutre.
\end{definition}

\begin{proposition}[R\'eduction de variance]
Les techniques principales sont :
\begin{enumerate}
    \item \textbf{Variables antith\'etiques} : pour chaque trajectoire
    $W_t$, utiliser aussi $-W_t$.
    \item \textbf{Variable de contr\^ole} : utiliser le prix connu de la
    g\'eom\'etrique pour corriger l'estimateur de l'arithm\'etique.
    \item \textbf{Importance sampling} : modifier la mesure pour \'echantillonner
    les \'ev\'enements rares.
\end{enumerate}
\end{proposition}

\begin{formules}[title=\textbf{Simulation de trajectoires sous $\Q$}]
Dans le mod\`ele de Black-Scholes :
\[
    S_{t_{i+1}} = S_{t_i}\exp\!\left[\left(r - \frac{\sigma^2}{2}\right)\Delta t
    + \sigma\sqrt{\Delta t}\,Z_i\right], \quad Z_i \sim \mathcal{N}(0, 1).
\]
\end{formules}

\section{M\'ethodes d'arbre pour les exotiques}

\begin{definition}[Arbre trinomial]
\index{arbre trinomial}
Un \textbf{arbre trinomial} discr\'etise le sous-jacent avec trois mouvements
\`a chaque pas : $S \to Su$, $S \to Sm$, $S \to Sd$. Les probabilit\'es
risque-neutres $(p_u, p_m, p_d)$ sont calibr\'ees pour reproduire les deux
premiers moments de $\ln S$.
\end{definition}

\begin{proposition}[Adaptation aux barri\`eres]
Pour une option barri\`ere, on ajuste l'arbre pour qu'un n\oe ud co\"incide
exactement avec la barri\`ere $B$, \'eliminant les erreurs de sp\'ecification
dues \`a la discr\'etisation.
\end{proposition}

\begin{codeblock}[title=\textbf{Prix d'une option asiatique par Monte Carlo}]
\begin{minted}{python}
import numpy as np

S0, K, r, sigma, T = 100, 100, 0.05, 0.2, 1.0
n_steps, n_paths = 252, 100_000
dt = T / n_steps

np.random.seed(42)
Z = np.random.randn(n_paths, n_steps)
S = np.zeros((n_paths, n_steps + 1))
S[:, 0] = S0

for i in range(n_steps):
    S[:, i+1] = S[:, i] * np.exp((r - 0.5*sigma**2)*dt
                                  + sigma*np.sqrt(dt)*Z[:, i])

S_mean = S[:, 1:].mean(axis=1)
payoff = np.maximum(S_mean - K, 0)
price = np.exp(-r*T) * payoff.mean()
stderr = np.exp(-r*T) * payoff.std() / np.sqrt(n_paths)

print(f"Prix asiatique (call): {price:.4f} +/- {1.96*stderr:.4f}")
\end{minted}
\end{codeblock}

\section{Exercices}

\begin{exercice}[Relation in-out]
D\'emontrer la relation in-out pour les options \`a barri\`ere en utilisant
l'absence d'opportunit\'e d'arbitrage.
\end{exercice}

\begin{exercice}[Asiatique vs vanille]
Montrer analytiquement que le prix d'un call asiatique (average price) est
inf\'erieur au prix du call vanille correspondant.
\end{exercice}

\begin{exercice}[Lookback pricing]
Impl\'ementer un pricer Monte Carlo pour une option lookback floating strike
et comparer avec la formule ferm\'ee de Goldman-Sosin-Gatto.
\end{exercice}

\begin{exercice}[R\'eplication d'une digitale]
Montrer qu'un call binaire cash-or-nothing de payoff $Q\,\ind_{S_T > K}$ peut
\^etre approch\'e par un call spread : $(C(K) - C(K + \epsilon))/\epsilon$
multipli\'e par $Q$, et calculer la limite $\epsilon \to 0$.
\end{exercice}

\begin{exercice}[Variables antith\'etiques]
Montrer que l'estimateur Monte Carlo avec variables antith\'etiques a une
variance inf\'erieure \`a l'estimateur standard si et seulement si les payoffs
sont n\'egativement corr\'el\'es sous l'inversion $W \to -W$.
\end{exercice}

\begin{exercice}[Barri\`ere discr\`ete]
Pour un down-and-out call avec barri\`ere discr\`ete observ\'ee
quotidiennement, appliquer la correction de Broadie-Glasserman-Kou :
$B_{\mathrm{eff}} = B\,e^{-\beta\sigma\sqrt{\Delta t}}$ avec
$\beta \approx 0{,}5826$ et comparer avec la barri\`ere continue.
\end{exercice}
